% =====================================================================
%  related_work_table.tex
%  Landscape of accelerators for structured sparse GEMV.
%
%  Feature matrix comparing this thesis (runtime-reconfigurable 2:M
%  sparse GEMV on an Alveo U280) against the eight-paper comparison set.
%  Every cell is backed by .claude/memory/RELATED_WORK.md, which records
%  the paper section each claim comes from (Lens 1-12).
%
%  USAGE
%    % =====================================================================
%  related_work_table.tex
%  Landscape of accelerators for structured sparse GEMV.
%
%  Feature matrix comparing this thesis (runtime-reconfigurable 2:M
%  sparse GEMV on an Alveo U280) against the eight-paper comparison set.
%  Every cell is backed by .claude/memory/RELATED_WORK.md, which records
%  the paper section each claim comes from (Lens 1-12).
%
%  USAGE
%    % =====================================================================
%  related_work_table.tex
%  Landscape of accelerators for structured sparse GEMV.
%
%  Feature matrix comparing this thesis (runtime-reconfigurable 2:M
%  sparse GEMV on an Alveo U280) against the eight-paper comparison set.
%  Every cell is backed by .claude/memory/RELATED_WORK.md, which records
%  the paper section each claim comes from (Lens 1-12).
%
%  USAGE
%    % =====================================================================
%  related_work_table.tex
%  Landscape of accelerators for structured sparse GEMV.
%
%  Feature matrix comparing this thesis (runtime-reconfigurable 2:M
%  sparse GEMV on an Alveo U280) against the eight-paper comparison set.
%  Every cell is backed by .claude/memory/RELATED_WORK.md, which records
%  the paper section each claim comes from (Lens 1-12).
%
%  USAGE
%    \input{Latex/related_work_table}
%
%  REQUIRED IN YOUR PREAMBLE
%    \usepackage{pifont}     % \ding{51} \ding{55}
%    \usepackage{xcolor}
%    \usepackage{booktabs}
%    \usepackage{multirow}
%
%  NOTE ON THE "(lim.)" MARKS
%    FNM-Trans and DynaX are marked \cmark(lim.), NOT \xmark, on the
%    runtime-switch column, and this is deliberate:
%      - FNM-Trans claims operation "within a single execution" (Sec. 1)
%        and its attention path is unambiguously runtime (Fig. 4), but
%        the paper never names a control signal -- the mode is implicit
%        in a self-describing binary-tree index (Sec. 4.2).
%      - DynaX varies X at runtime via thresholds (Eq. 2), but N:M is
%        fixed at synthesis (Sec. 6.1).
%    A bare \xmark on either is an overclaim an examiner who has read
%    the papers will catch, and it is not needed: both already fail
%    four other columns.
%
%  NOTE ON WIDTH
%    If seven columns is too wide, drop the first ("Struct. N:M
%    Weights"). It is the "is it in the family" filter; the remaining
%    six still leave every other row with at least one cross.
% =====================================================================

\providecommand{\cmark}{\textcolor{green!60!black}{\ding{51}}}
\providecommand{\xmark}{\textcolor{red}{\ding{55}}}

\begin{table}[t]
\centering
\small
\setlength{\tabcolsep}{4pt}
\begin{tabular}{lccccccc}
\toprule
\multirow{3}{*}{Architecture} & Struct.  & Runtime         & Group             & HBM     & No      & Real Si.\ + & Iso-Platform \\
                              & $N{:}M$  & Ratio           & $M\!\ge\!32$      & Multi-  & Reduct. & Meas.       & Dense        \\
                              & Weights  & Switch          & (93.75\%)         & Channel & Tree    & Power       & Baseline     \\
\midrule
Mishra~\cite{mishra2021}        & \cmark & \xmark          & \xmark & \xmark & \cmark & \xmark & \cmark \\
FNM-Trans~\cite{fnmtrans2024}   & \cmark & \cmark\,(lim.)  & \xmark & \xmark & \xmark & \xmark & \xmark \\
Serpens~\cite{serpens2022}      & \xmark & \xmark          & \xmark & \cmark & \xmark & \cmark & \xmark \\
SST~\cite{sst2025}              & \cmark & \cmark          & \xmark & \xmark & \xmark & \xmark & \xmark \\
Transitive Array~\cite{ta2025}  & \xmark & \xmark          & \xmark & \xmark & \xmark & \xmark & \xmark \\
DFX~\cite{dfx2022}              & \xmark & \xmark          & \xmark & \cmark & \xmark & \cmark & \xmark \\
DynaX~\cite{dynax2025}          & \cmark & \cmark\,(lim.)  & \cmark & \xmark & \xmark & \xmark & \xmark \\
FlightVGM~\cite{flightvgm2025}  & \xmark & \cmark          & \xmark & \cmark & \xmark & \cmark & \xmark \\
\midrule
\textbf{This work}              & \cmark & \cmark          & \cmark & \cmark & \cmark & \cmark & \cmark \\
\bottomrule
\end{tabular}
\caption{Landscape of accelerators for structured sparse GEMV.
\emph{Struct.\ $N{:}M$ Weights}: a fixed, hardware-friendly sparsity
pattern imposed on the weight tensor.
\emph{Runtime Ratio Switch}: the sparsity ratio changes without
re-synthesis.
\emph{Group $M\!\ge\!32$}: the group size reaches 32, i.e.\ a 93.75\%
sparsity ceiling rather than the 75\% of $M\!=\!4$ designs.
\emph{HBM Multi-Channel}: operands streamed over HBM pseudo-channels.
\emph{No Reduction Tree}: load balance by construction, with no
arbiters, schedulers, crossbars or adder trees.
\emph{Real Si.\ + Meas.\ Power}: deployed on silicon with power
measured in-system rather than modelled or read from a synthesis
report.
\emph{Iso-Platform Dense Baseline}: sparse compared against a dense
counterpart on the same board, flow and clock.}
\label{tab:related-work}
\end{table}

%
%  REQUIRED IN YOUR PREAMBLE
%    \usepackage{pifont}     % \ding{51} \ding{55}
%    \usepackage{xcolor}
%    \usepackage{booktabs}
%    \usepackage{multirow}
%
%  NOTE ON THE "(lim.)" MARKS
%    FNM-Trans and DynaX are marked \cmark(lim.), NOT \xmark, on the
%    runtime-switch column, and this is deliberate:
%      - FNM-Trans claims operation "within a single execution" (Sec. 1)
%        and its attention path is unambiguously runtime (Fig. 4), but
%        the paper never names a control signal -- the mode is implicit
%        in a self-describing binary-tree index (Sec. 4.2).
%      - DynaX varies X at runtime via thresholds (Eq. 2), but N:M is
%        fixed at synthesis (Sec. 6.1).
%    A bare \xmark on either is an overclaim an examiner who has read
%    the papers will catch, and it is not needed: both already fail
%    four other columns.
%
%  NOTE ON WIDTH
%    If seven columns is too wide, drop the first ("Struct. N:M
%    Weights"). It is the "is it in the family" filter; the remaining
%    six still leave every other row with at least one cross.
% =====================================================================

\providecommand{\cmark}{\textcolor{green!60!black}{\ding{51}}}
\providecommand{\xmark}{\textcolor{red}{\ding{55}}}

\begin{table}[t]
\centering
\small
\setlength{\tabcolsep}{4pt}
\begin{tabular}{lccccccc}
\toprule
\multirow{3}{*}{Architecture} & Struct.  & Runtime         & Group             & HBM     & No      & Real Si.\ + & Iso-Platform \\
                              & $N{:}M$  & Ratio           & $M\!\ge\!32$      & Multi-  & Reduct. & Meas.       & Dense        \\
                              & Weights  & Switch          & (93.75\%)         & Channel & Tree    & Power       & Baseline     \\
\midrule
Mishra~\cite{mishra2021}        & \cmark & \xmark          & \xmark & \xmark & \cmark & \xmark & \cmark \\
FNM-Trans~\cite{fnmtrans2024}   & \cmark & \cmark\,(lim.)  & \xmark & \xmark & \xmark & \xmark & \xmark \\
Serpens~\cite{serpens2022}      & \xmark & \xmark          & \xmark & \cmark & \xmark & \cmark & \xmark \\
SST~\cite{sst2025}              & \cmark & \cmark          & \xmark & \xmark & \xmark & \xmark & \xmark \\
Transitive Array~\cite{ta2025}  & \xmark & \xmark          & \xmark & \xmark & \xmark & \xmark & \xmark \\
DFX~\cite{dfx2022}              & \xmark & \xmark          & \xmark & \cmark & \xmark & \cmark & \xmark \\
DynaX~\cite{dynax2025}          & \cmark & \cmark\,(lim.)  & \cmark & \xmark & \xmark & \xmark & \xmark \\
FlightVGM~\cite{flightvgm2025}  & \xmark & \cmark          & \xmark & \cmark & \xmark & \cmark & \xmark \\
\midrule
\textbf{This work}              & \cmark & \cmark          & \cmark & \cmark & \cmark & \cmark & \cmark \\
\bottomrule
\end{tabular}
\caption{Landscape of accelerators for structured sparse GEMV.
\emph{Struct.\ $N{:}M$ Weights}: a fixed, hardware-friendly sparsity
pattern imposed on the weight tensor.
\emph{Runtime Ratio Switch}: the sparsity ratio changes without
re-synthesis.
\emph{Group $M\!\ge\!32$}: the group size reaches 32, i.e.\ a 93.75\%
sparsity ceiling rather than the 75\% of $M\!=\!4$ designs.
\emph{HBM Multi-Channel}: operands streamed over HBM pseudo-channels.
\emph{No Reduction Tree}: load balance by construction, with no
arbiters, schedulers, crossbars or adder trees.
\emph{Real Si.\ + Meas.\ Power}: deployed on silicon with power
measured in-system rather than modelled or read from a synthesis
report.
\emph{Iso-Platform Dense Baseline}: sparse compared against a dense
counterpart on the same board, flow and clock.}
\label{tab:related-work}
\end{table}

%
%  REQUIRED IN YOUR PREAMBLE
%    \usepackage{pifont}     % \ding{51} \ding{55}
%    \usepackage{xcolor}
%    \usepackage{booktabs}
%    \usepackage{multirow}
%
%  NOTE ON THE "(lim.)" MARKS
%    FNM-Trans and DynaX are marked \cmark(lim.), NOT \xmark, on the
%    runtime-switch column, and this is deliberate:
%      - FNM-Trans claims operation "within a single execution" (Sec. 1)
%        and its attention path is unambiguously runtime (Fig. 4), but
%        the paper never names a control signal -- the mode is implicit
%        in a self-describing binary-tree index (Sec. 4.2).
%      - DynaX varies X at runtime via thresholds (Eq. 2), but N:M is
%        fixed at synthesis (Sec. 6.1).
%    A bare \xmark on either is an overclaim an examiner who has read
%    the papers will catch, and it is not needed: both already fail
%    four other columns.
%
%  NOTE ON WIDTH
%    If seven columns is too wide, drop the first ("Struct. N:M
%    Weights"). It is the "is it in the family" filter; the remaining
%    six still leave every other row with at least one cross.
% =====================================================================

\providecommand{\cmark}{\textcolor{green!60!black}{\ding{51}}}
\providecommand{\xmark}{\textcolor{red}{\ding{55}}}

\begin{table}[t]
\centering
\small
\setlength{\tabcolsep}{4pt}
\begin{tabular}{lccccccc}
\toprule
\multirow{3}{*}{Architecture} & Struct.  & Runtime         & Group             & HBM     & No      & Real Si.\ + & Iso-Platform \\
                              & $N{:}M$  & Ratio           & $M\!\ge\!32$      & Multi-  & Reduct. & Meas.       & Dense        \\
                              & Weights  & Switch          & (93.75\%)         & Channel & Tree    & Power       & Baseline     \\
\midrule
Mishra~\cite{mishra2021}        & \cmark & \xmark          & \xmark & \xmark & \cmark & \xmark & \cmark \\
FNM-Trans~\cite{fnmtrans2024}   & \cmark & \cmark\,(lim.)  & \xmark & \xmark & \xmark & \xmark & \xmark \\
Serpens~\cite{serpens2022}      & \xmark & \xmark          & \xmark & \cmark & \xmark & \cmark & \xmark \\
SST~\cite{sst2025}              & \cmark & \cmark          & \xmark & \xmark & \xmark & \xmark & \xmark \\
Transitive Array~\cite{ta2025}  & \xmark & \xmark          & \xmark & \xmark & \xmark & \xmark & \xmark \\
DFX~\cite{dfx2022}              & \xmark & \xmark          & \xmark & \cmark & \xmark & \cmark & \xmark \\
DynaX~\cite{dynax2025}          & \cmark & \cmark\,(lim.)  & \cmark & \xmark & \xmark & \xmark & \xmark \\
FlightVGM~\cite{flightvgm2025}  & \xmark & \cmark          & \xmark & \cmark & \xmark & \cmark & \xmark \\
\midrule
\textbf{This work}              & \cmark & \cmark          & \cmark & \cmark & \cmark & \cmark & \cmark \\
\bottomrule
\end{tabular}
\caption{Landscape of accelerators for structured sparse GEMV.
\emph{Struct.\ $N{:}M$ Weights}: a fixed, hardware-friendly sparsity
pattern imposed on the weight tensor.
\emph{Runtime Ratio Switch}: the sparsity ratio changes without
re-synthesis.
\emph{Group $M\!\ge\!32$}: the group size reaches 32, i.e.\ a 93.75\%
sparsity ceiling rather than the 75\% of $M\!=\!4$ designs.
\emph{HBM Multi-Channel}: operands streamed over HBM pseudo-channels.
\emph{No Reduction Tree}: load balance by construction, with no
arbiters, schedulers, crossbars or adder trees.
\emph{Real Si.\ + Meas.\ Power}: deployed on silicon with power
measured in-system rather than modelled or read from a synthesis
report.
\emph{Iso-Platform Dense Baseline}: sparse compared against a dense
counterpart on the same board, flow and clock.}
\label{tab:related-work}
\end{table}

%
%  REQUIRED IN YOUR PREAMBLE
%    \usepackage{pifont}     % \ding{51} \ding{55}
%    \usepackage{xcolor}
%    \usepackage{booktabs}
%    \usepackage{multirow}
%
%  NOTE ON THE "(lim.)" MARKS
%    FNM-Trans and DynaX are marked \cmark(lim.), NOT \xmark, on the
%    runtime-switch column, and this is deliberate:
%      - FNM-Trans claims operation "within a single execution" (Sec. 1)
%        and its attention path is unambiguously runtime (Fig. 4), but
%        the paper never names a control signal -- the mode is implicit
%        in a self-describing binary-tree index (Sec. 4.2).
%      - DynaX varies X at runtime via thresholds (Eq. 2), but N:M is
%        fixed at synthesis (Sec. 6.1).
%    A bare \xmark on either is an overclaim an examiner who has read
%    the papers will catch, and it is not needed: both already fail
%    four other columns.
%
%  NOTE ON WIDTH
%    If seven columns is too wide, drop the first ("Struct. N:M
%    Weights"). It is the "is it in the family" filter; the remaining
%    six still leave every other row with at least one cross.
% =====================================================================

\providecommand{\cmark}{\textcolor{green!60!black}{\ding{51}}}
\providecommand{\xmark}{\textcolor{red}{\ding{55}}}

\begin{table}[t]
\centering
\small
\setlength{\tabcolsep}{4pt}
\begin{tabular}{lccccccc}
\toprule
\multirow{3}{*}{Architecture} & Struct.  & Runtime         & Group             & HBM     & No      & Real Si.\ + & Iso-Platform \\
                              & $N{:}M$  & Ratio           & $M\!\ge\!32$      & Multi-  & Reduct. & Meas.       & Dense        \\
                              & Weights  & Switch          & (93.75\%)         & Channel & Tree    & Power       & Baseline     \\
\midrule
Mishra~\cite{mishra2021}        & \cmark & \xmark          & \xmark & \xmark & \cmark & \xmark & \cmark \\
FNM-Trans~\cite{fnmtrans2024}   & \cmark & \cmark\,(lim.)  & \xmark & \xmark & \xmark & \xmark & \xmark \\
Serpens~\cite{serpens2022}      & \xmark & \xmark          & \xmark & \cmark & \xmark & \cmark & \xmark \\
SST~\cite{sst2025}              & \cmark & \cmark          & \xmark & \xmark & \xmark & \xmark & \xmark \\
Transitive Array~\cite{ta2025}  & \xmark & \xmark          & \xmark & \xmark & \xmark & \xmark & \xmark \\
DFX~\cite{dfx2022}              & \xmark & \xmark          & \xmark & \cmark & \xmark & \cmark & \xmark \\
DynaX~\cite{dynax2025}          & \cmark & \cmark\,(lim.)  & \cmark & \xmark & \xmark & \xmark & \xmark \\
FlightVGM~\cite{flightvgm2025}  & \xmark & \cmark          & \xmark & \cmark & \xmark & \cmark & \xmark \\
\midrule
\textbf{This work}              & \cmark & \cmark          & \cmark & \cmark & \cmark & \cmark & \cmark \\
\bottomrule
\end{tabular}
\caption{Landscape of accelerators for structured sparse GEMV.
\emph{Struct.\ $N{:}M$ Weights}: a fixed, hardware-friendly sparsity
pattern imposed on the weight tensor.
\emph{Runtime Ratio Switch}: the sparsity ratio changes without
re-synthesis.
\emph{Group $M\!\ge\!32$}: the group size reaches 32, i.e.\ a 93.75\%
sparsity ceiling rather than the 75\% of $M\!=\!4$ designs.
\emph{HBM Multi-Channel}: operands streamed over HBM pseudo-channels.
\emph{No Reduction Tree}: load balance by construction, with no
arbiters, schedulers, crossbars or adder trees.
\emph{Real Si.\ + Meas.\ Power}: deployed on silicon with power
measured in-system rather than modelled or read from a synthesis
report.
\emph{Iso-Platform Dense Baseline}: sparse compared against a dense
counterpart on the same board, flow and clock.}
\label{tab:related-work}
\end{table}
